
\section{问题分析}

三问都要求安排2024至2030年的种植面积，并共享地块容量、适种性、重茬、豆类覆盖和水浇地模式等约束。它们的区别主要在经济参数的确定方式和结果评价范围。问题一处理稳定参数下的两种销售规则，问题二引入产销参数波动，问题三进一步考察作物关系和经济变量相关性。

模型以逐年逐季逐地块的种植面积为统一变量。求解后独立核验容量、适种、重茬、豆类覆盖和水浇地模式，评价环节计算累计净收益、5\%分位数、最差5\%平均收益和亏损情景比例。每一问都保留能完成同一任务的基线，并根据相同指标进行比较。

附件提供54个地块、41种作物、87条2023年种植记录和107条统计参数；价格取区间中点作为确定性代表值，预期销量以2023年种植面积与对应亩产量汇总构造，智慧大棚第一季缺少的参数由普通大棚同季参数补齐。

\subsection{问题一分析}

问题一的输出是两套完整面积表，分别对应超额产量滞销和半价销售。核心决策是在七年范围内分配各地块、各季次的作物面积，使累计净收益尽可能高，同时满足题目规定的土地利用和轮作要求。由于同一地块当前季的选择会影响下一季重茬限制，豆类覆盖又跨越连续三年，逐年独立安排不能保证整体可行。

本问的关键在于正确表达分段销售收益和跨期约束。产量不超过销量上限时按正常价格销售，超出部分的价值由销售模式决定；同一面积方案在两种规则下可能形成不同的作物集中程度。结果验收不仅比较累计净收益，还需逐格检查面积容量、作物适宜性、相邻季重茬、豆类窗口和水浇地种植模式，并报告未闭合的 MIP Gap。

\subsection{问题二分析}

问题二仍需输出唯一一套2024至2030年面积表，但销量、亩产量、成本和价格会随年份变化。题目给出了变化区间或趋势，却没有提供概率分布和多年样本。因此，不能把某一条随机路径当成已知未来，也不能让种植面积随测试情景事后调整。

本问需要区分方案生成和风险评价。先依据题目区间形成透明的中心趋势，用于求解固定方案；再把该方案放入共同的样本外情景，与保持2023年参数不变的基线配对比较。除平均累计收益外，还要检查5\%分位数、最差5\%平均收益、最低收益和亏损情景比例。由此得到的结论仅适用于设定的波动范围，均值优势不能自动解释为下尾风险更低。

\subsection{问题三分析}

问题三在问题二基础上增加作物替代、互补关系以及销量、亩产量和价格之间的相关性。附件没有提供可直接估计这些关系的历史样本，因此本问不是识别真实市场弹性，而是在明确的模拟关系下考察种植方案和收益分布如何变化。

为分离新增关系结构的影响，问题三保持问题二的土地约束、管理配置、销售规则和风险指标不变，并用关闭关系结构的问题二方案作为基线。比较内容包括平均收益、下尾风险和种植面积变化。逐地块重合度容易受到多组等价解影响，因此还需比较全乡村41种作物的七年累计面积结构。只有收益和宏观面积结构同时发生明显变化，才能将差异解释为结构性跳变；否则更适合视为地块层面的调度调整。
